% !PRINTED none
\begin{center}
\gesperrt{Vierter Abschnitt.}
\end{center}

\begin{center}
\gesperrt{Cubische und biquadratische Abel'sche K\"orper.}
\end{center}

\begin{center}
\rule{3cm}{0.4pt}
\end{center}

\begin{center}
\S.\,23.
\end{center}

\begin{center}
\gesperrt{Cubische Kreistheilungsk\"orper.}
\end{center}

Die allgemeinen Untersuchungen, die in den letzten Paragraphen dargestellt sind, gestatten mannigfache Anwendungen auf specielle Fälle, von denen die einfachsten hier betrachtet werden sollen. Es ist dabei bemerkenswerth, dass schon die einfachsten Fälle, wenn man sie direct angreift, fast in dieselben Schwierigkeiten hineinführen, die wir durch die allgemeine Untersuchung überwunden haben.

Wir wollen zuerst die Aufgabe behandeln, alle \gesperrt{cubischen Kreistheilungskörper, also alle Kreistheilungsperioden, die einer rationalen Gleichung $3^{\mathrm{ten}}$ Grades genügen,} aufzufinden, und zwar so, dass von mehreren Ausdrücken, deren dieselbe Grösse fähig ist, der einfachste gesetzt wird.

Wir haben hier nur eine Invariante $i_1=3$, und nach den Vorschriften in der Zusammenstellung des \S.\,21 dürfen in $m$ aufgenommen werden der Factor 9, nicht aber 3, ferner Primzahlen in beliebiger Menge von der Form $6N+1$, also die Primzahlen 7, 13, 19, 31, 37, 43, 61, 67, 73, 79, 97 . . ., und so bekommen wir also als geeignete Werthe von $m$ unter 200:
\[
\begin{aligned}
m &= 7,\ 9,\ 13,\ 19,\ 31,\ 37,\ 43,\ 61,\ \mathbf{63},\ 67,\ 73,\ 79,\ \mathbf{91},\ 97,\ 103,\ 109,\\
&\quad \mathbf{117},\ 127,\ \mathbf{133},\ 139,\ 151,\ 157,\ 163,\ \mathbf{171},\ 181,\ 193,\ 199,
\end{aligned}
\]
unter denen die Zahlen mit mehreren Primtheilern durch fetten Druck ausgezeichnet sind. Die kleinste Zahl $m$, in der mehr als zwei Primzahlen aufgehen, ist $7.9.13=819$.
